% numerals-ok: begin
\clearpage
\section{Sensitivity of the two decision layers}
\label{app:sensitivity}

Two parameters carry the decision chain and neither is estimated inside the paper: the
service multiplier that maps a duration into a service time, and the decoding constants
of the boundary detector. The tables below move each of them and report what changes.

\begin{table}[H]
\centering
\scriptsize
\begin{tabular}{@{}lrrrrr@{}}
\hline
Case & $\alpha$ & jobs & mean predicted & realised & gap \\
\hline
S1 & 0.75 & 720 & 0.9701 & 1.0000 & -0.0299 \\
S1 & 1.00 & 720 & 0.9290 & 0.9750 & -0.0460 \\
S1 & 1.50 & 628 & 0.8024 & 0.8344 & -0.0320 \\
S2 & 0.75 & 1062 & 0.9179 & 0.9793 & -0.0614 \\
S2 & 1.00 & 903 & 0.8067 & 0.8538 & -0.0471 \\
S2 & 1.50 & 665 & 0.7411 & 0.7684 & -0.0273 \\
S3 & 0.75 & 1315 & 0.9179 & 0.9567 & -0.0388 \\
S3 & 1.00 & 1148 & 0.7652 & 0.8118 & -0.0467 \\
S3 & 1.50 & 839 & 0.7040 & 0.7318 & -0.0278 \\
\hline
\end{tabular}
\caption{Service-multiplier sweep on the closed-form layer. Each cell re-prices the same
executed jobs at $\alpha\in\{0.75,1.0,1.5\}$; the largest absolute gap is
0.0614 and the mean absolute gap 0.0397.}
\label{tab:alpha-sweep}
\end{table}

\begin{table}[H]
\centering
\scriptsize
\setlength{\tabcolsep}{4pt}
\begin{tabular}{@{}llrrrr@{}}
\hline
Family & setting & $\sigma$ & $D$ & validation ARI & train ARI \\
\hline
quant & q0.85 & 3 & 78 & 0.6010 & 0.7147 \\
fixed & H0.20 & 14 & 78 & 0.5969 & 0.6909 \\
fixed & H0.25 & 7 & 78 & 0.5939 & 0.7175 \\
quant & q0.85 & 5 & 78 & 0.5918 & 0.7029 \\
fixed & H0.25 & 2 & 78 & 0.5900 & 0.6984 \\
quant & q0.90 & 3 & 78 & 0.5880 & 0.7333 \\
quant & q0.90 & 2 & 78 & 0.5875 & 0.7256 \\
fixed & H0.20 & 10 & 78 & 0.5874 & 0.6882 \\
fixed & H0.30 & 5 & 78 & 0.5870 & 0.7401 \\
greedy & lam0.30 & 5 & 78 & 0.5870 & 0.7401 \\
\hline
\end{tabular}
\caption{The ten best of the 192 decoder configurations on validation records. The
configuration that would be best on test ranks 24 here, which is the transfer cost we discuss under Problem~1.}
\label{tab:decoder-sweep}
\end{table}

The two sweeps answer different questions. The first asks whether the closed form is
calibrated when the service multiplier is wrong: it stays within 0.0614 in the worst cell and the order of cases
never changes. The second asks whether the decoding constants transfer: they do not,
and the gap is the price of choosing them on seven validation records.


\subsection*{The reservation is not a third parameter}

Figure~\ref{fig:policy} first fixes the policy set; Figure~\ref{fig:phase} then sweeps the two axes a practitioner would tune: the service
multiplier $\alpha$ and the residual quantile $q$ that sets the reserve. Raising $q$
from 0.30 to 0.90 costs \slotPPhaseDropSTwo{} points in S\slotCaseTwo{} and
\slotPPhaseDropSThree{} points in S\slotCaseThree{} at fixed $\alpha$; moving $\alpha$
from 0.75 to 1.50 moves the rate by \slotPPhaseAlphaSpreadSTwo{} and
\slotPPhaseAlphaSpreadSThree{} points. The rate is flat in $q$ while the reserve stays
below the realised service time and falls once it exceeds it, which is the empirical
form of Proposition~3: the reservation is a liability, not a tuning knob.

\clearpage
\subsection*{The value ladder in one picture}

Figure~\ref{fig:value} puts the four mechanisms on one axis, in the same units.
The scheduling rule is worth \slotPValueRuleSTwo{} points in S\slotCaseTwo{} and
\slotPValueRuleSThree{} in S\slotCaseThree{}; handing the same rule the true duration is
worth \slotPValueTruthSTwo{} and \slotPValueTruthSThree{}; the paid review is worth
nothing detectable (best cell +0.28 point, interval crossing zero); and raising the
reservation quantile from 0.30 to 0.90 moves the rate by
\slotPValueReserveSTwo{} and \slotPValueReserveSThree{} points. The ordering is the
paper in one line: rule, then knowledge, then nothing, then a liability.

% numerals-ok: end
